\section{Weak second order, infinity, and the admissible fragment}
\label{sec:p03-hyp-characterization}

Write \(L_{\mathrm{HYP}}\) for the infinitary formulas whose codes
belong to \(\mathrm{HYP}\), restricted here to finite signatures.
Symbols are coded by natural numbers with their kinds and arities.
For any other finite naming of a signature, transport the syntax along
an arity-preserving renaming to such codes. Two choices differ by a
finite coded renaming, which preserves membership in \(\mathrm{HYP}\)
by \cref{thm:p03-adm-syntax}; thus expressive power does not depend on
the choice of names.
Write \(L_w\) for weak second-order logic, with quantifiers over
finite subsets, and \(L(Q_0)\) for first-order logic with the
infinitely-many quantifier. The following equality is an equality of
expressive power on each finite signature, not an identification of
the three syntactic alphabets.

\begin{hblemma}{thm:p03-hyp-projective}{proved=annex:p3.hyp.1}{Projective presentation of admissible truth}
For every sentence \(\varphi\in L_{\mathrm{HYP}}\), its model class
and its complement have finite-signature relativized projective
presentations in each of \(L_w\) and \(L(Q_0)\).
The construction allows an empty model domain when its signature permits
one. Its auxiliary set-theoretic model always has a nonempty domain.
\end{hblemma}

\begin{hbtheorem}{thm:p03-hyp-characterization}{proved=annex:p3.hyp.2}{The admissible closure characterization}
On finite signatures, with relativized projective witnesses,
\[
 \Delta(L_w)=\Delta(L(Q_0))=L_{\mathrm{HYP}}.
\]
Thus \(L_{\mathrm{HYP}}\) is the least Souslin--Kleene extension of
either logic. Its admissible interpolation and ordinary Beth properties
are the separate results of
\cref{thm:p03-adm-interpolation,thm:p03-adm-beth}; the leastness assertion
here concerns the Souslin--Kleene property.
\end{hbtheorem}

The target characterization is Barwise's Theorem~4.1
\cite[Part~II, \S4, pp.~260--262]{Barwise1974Axioms}.
The proof here supplies its required fixed names, uniform satisfaction,
and admissible separation explicitly. For finite auxiliary presentation
it uses the external truth relation and standard-number construction of
\cref{thm:p03-recursive-auxiliary-presentation}. This is a stated
alternative to the source's single pure set-theoretic truncation
sentence, not a claim that those two finite presentations are identical.
\(\mathrm{HYP}\) and the ambient models of set theory are used to
prove the equality; they are not new primitives of either of the two
input logics.
